\documentclass[11pt]{article}

\usepackage[margin=1in]{geometry}
\usepackage{amsmath,amssymb,amsthm,mathtools}
\usepackage{graphicx}
\usepackage{microtype}
\usepackage[hidelinks]{hyperref}
\usepackage{enumitem}

\newtheorem{theorem}{Theorem}
\newtheorem{lemma}{Lemma}
\newtheorem{remark}{Remark}
\newcommand{\C}{\mathbb C}
\newcommand{\Mtwo}{M_2(\C)}
\newcommand{\Ep}{E_p}
\newcommand{\dist}{\operatorname{dist}}
\newcommand{\Tr}{\operatorname{Tr}}
\newcommand{\diag}{\operatorname{diag}}

\title{Three-Root Algebraic Components at Distance One Half}
\author{Candidate counterexample for arXiv:1601.01505\\
Run \texttt{fa\_banach\_001}, lane 0}
\date{9 August 2026}

\begin{document}
\maketitle

\begin{center}
\fbox{\parbox{0.91\textwidth}{\textbf{Status.}
Candidate counterexample, likely valid, pending specialist human review.
For $p(\lambda)=\lambda(\lambda-1)(\lambda-2)$ in $M_2(\C)$, two distinct
components of $E_p(M_2(\C))$ have distance exactly $1/2$, while the source
conjectures the lower bound $1$.}}
\end{center}

\section{Source conjecture}

For a unital complex Banach algebra $A$ and a polynomial
\[
 p(\lambda)=\prod_{j=1}^n(\lambda-\lambda_j)
\]
with distinct roots, Makai and Zem\'anek write
\[
 E_p(A)=\{a\in A:p(a)=0\}.
\]
They conjecture that any two distinct connected components $C_1,C_2$ of
$E_p(A)$ satisfy
\begin{equation}\label{eq:source}
 \dist(C_1,C_2)\geq
 \min_{i\neq j}|\lambda_i-\lambda_j|.
\end{equation}
They state the analogous conjecture for self-adjoint algebraic elements in a
$C^*$-algebra \cite[p.~5]{MZ16}.

\begin{figure}[ht]
  \centering
  \includegraphics[width=\textwidth]{figures/open_problem_crop.png}
  \caption{Theorem 9 and the conjecture in arXiv:1601.01505, PDF page 5.}
\end{figure}

\section{Proof intuition}

With three roots, two components can have spectral multiplicities whose traces
differ by only one even though their occupied nonzero roots differ.  We choose
rank-one idempotents $e,f$ so that the difference between the root-$1$ element
$e$ and the root-$2$ element $2f$ has the form
\[
 2f-e=\frac12 I+N,
\]
where $N$ is nonzero but square-zero.  A nonunitary similarity can shrink a
$2\times2$ nilpotent Jordan block by an arbitrary factor.  Simultaneously
conjugating $e$ and $2f$ preserves their connected components, while their
difference tends to $I/2$.  The trace difference supplies the matching lower
bound $1/2$.

\section{Counterexample}

Equip $A=\Mtwo$ with its usual operator norm and take
\[
 p(\lambda)=\lambda(\lambda-1)(\lambda-2).
\]
Let $C_{01}$ denote the component of elements with root multiplicities
$(1,1,0)$ and let $C_{02}$ denote the component with multiplicities
$(1,0,1)$, in the root order $(0,1,2)$.

\begin{theorem}\label{thm:main}
The components $C_{01}$ and $C_{02}$ are distinct and
\[
 \dist(C_{01},C_{02})=\frac12.
\]
Since the minimum separation of the roots $0,1,2$ is $1$, this disproves
\eqref{eq:source} for the Banach-algebra set $E_p(A)$.
\end{theorem}

\begin{proof}
Set
\[
 e=\begin{pmatrix}1&0\\0&0\end{pmatrix},
 \qquad
 f=\begin{pmatrix}\frac98&1\\[2pt]-\frac9{64}&-\frac18\end{pmatrix}.
\]
A direct multiplication gives $e^2=e$ and $f^2=f$, and both idempotents have
rank one.  Thus $e\in E_p(A)$ has root multiplicities $(1,1,0)$, whereas
$2f\in E_p(A)$ has root multiplicities $(1,0,1)$.

For completeness, these multiplicities are constant on connected components.
Indeed, for $a\in E_p(A)$ the Lagrange spectral idempotents
\begin{equation}\label{eq:riesz}
 q_j(a)=\prod_{k\neq j}\frac{a-\lambda_kI}{\lambda_j-\lambda_k}
 \qquad(j=0,1,2)
\end{equation}
depend polynomially on $a$.  Their ranks are locally constant, hence constant
on every connected set.  Since the rank triples for $e$ and $2f$ differ, they
belong to distinct components.

Now calculate
\begin{align*}
 D:=2f-e
 &=\begin{pmatrix}\frac54&2\\[2pt]-\frac9{32}&-\frac14\end{pmatrix}
   =\frac12I+N,\\
 N&=\begin{pmatrix}\frac34&2\\[2pt]-\frac9{32}&-\frac34\end{pmatrix}.
\end{align*}
Then $N\neq0$ and $N^2=0$.  More explicitly, with
\[
 P=\begin{pmatrix}8&0\\-3&4\end{pmatrix},
 \qquad J=\begin{pmatrix}0&1\\0&0\end{pmatrix},
\]
we have $P^{-1}NP=J$.

For $t>0$, define
\[
 S_t=P\begin{pmatrix}t&0\\0&1\end{pmatrix}P^{-1}
 =\begin{pmatrix}t&0\\[2pt]\frac{3(1-t)}8&1\end{pmatrix}.
\]
Then $S_t$ is invertible and
\begin{equation}\label{eq:shrink}
 S_tNS_t^{-1}=tN.
\end{equation}
Put
\[
 a_t=S_teS_t^{-1},
 \qquad b_t=2S_tfS_t^{-1}.
\]
The group $GL_2(\C)$ is path connected, so similarity by $S_t$ connects
$a_t$ to $e$ inside $E_p(A)$ and connects $b_t$ to $2f$ inside $E_p(A)$.
Therefore $a_t\in C_{01}$ and $b_t\in C_{02}$.  Equations above give
\[
 b_t-a_t=S_tDS_t^{-1}=\frac12I+tN,
\]
and hence
\[
 \dist(C_{01},C_{02})leq
 \lim_{t\downarrow0}\|b_t-a_t\|=\frac12.
\]

For the reverse inequality, every $a\in C_{01}$ has ordinary matrix trace
$\Tr(a)=1$, while every $b\in C_{02}$ has $\Tr(b)=2$.  Since
\[
 |\Tr(x)|\leq 2\|x\|\qquad(x\in\Mtwo),
\]
we obtain
\[
 \|b-a\|\geq\frac{|\Tr(b-a)|}{2}=\frac12
 \qquad(a\in C_{01},\ b\in C_{02}).
\]
The two bounds prove the asserted exact distance.
\end{proof}

\begin{remark}[A fixed pair already violates the conjecture]
The square-zero matrix $N$ has operator norm $73/32$.  Taking $t=1/10$ in
the construction gives
\[
 \|b_{1/10}-a_{1/10}\|
 \leq \frac12+\frac1{10}\frac{73}{32}
 =\frac{233}{320}<1.
\]
Thus one need not pass to an infimum to exhibit a cross-component pair closer
than the minimum root separation.
\end{remark}

\section{Verification notes}

\begin{enumerate}[leftmargin=*]
\item \textbf{Exact algebraic identities.}  The idempotency of $e,f$, the
identity $N^2=0$, and the conjugation formula \eqref{eq:shrink} are rational
$2\times2$ matrix calculations.
\item \textbf{Component invariant.}  Formula \eqref{eq:riesz} avoids any
appeal to a delicate global classification: the ranks of the spectral
idempotents alone prove that the components are distinct.
\item \textbf{Component membership of the families.}  The paths come from
paths in $GL_2(\C)$ and therefore remain in $E_p(A)$ under similarity.
\item \textbf{Sharp lower bound.}  The trace estimate applies to every pair
from the two components, not only to the constructed families.
\item \textbf{Supplementary code.}  The script
\texttt{code/verify\_matrices.py} checks the rational identities symbolically
and samples the norm convergence.  It is not a dependency of the proof.
\end{enumerate}

\section{\hspace{0.25em}Scope and limitations}

The counterexample disproves the Banach-algebra $E_p(A)$ assertion, even with
the finite-dimensional $C^*$-algebra $A=M_2(\C)$ as ambient algebra.  The
elements $a_t,b_t$ are generally not self-adjoint because $S_t$ is not
unitary.  Consequently, this packet does \emph{not} disprove the separate
self-adjoint conjecture for $S_p(A)$.  It also does not decide the weaker
question, mentioned by the source authors, whether distinct Banach-algebra
components always have some positive separation depending on the roots.

\section{\hspace{0.25em}Novelty check}

A bounded search on 9 August 2026 covered:
\begin{itemize}[leftmargin=*]
\item the run's lightweight indexes and local 2000--2026 arXiv source corpus;
\item exact and close-variant web/arXiv searches using the source title,
authors, ``distance of connected components'', and the $E_p(A)$ notation;
\item the authors' expanded 2018 preprint, which repeats the same distance
problem \cite{MZ18};
\item the OpenAlex citation graph of DOI
\texttt{10.1007/s10587-016-0294-6}.  It contained two substantive citing
works, on idempotents and similarity-orbit commutativity; neither search result
addressed the component-distance conjecture.
\end{itemize}
No prior explicit counterexample or resolution was found.  This bounded search
does not prove novelty, and the candidate should be checked by a specialist.

\section{\hspace{0.25em}Human-review recommendation}

\textbf{Recommendation: high-priority verification.}  The construction is
finite-dimensional and the proof is short.  A reviewer should verify the four
matrix identities, confirm that the spectral-idempotent rank triples separate
the components, and check that simultaneous similarity produces paths within
the two fixed components.

\begin{thebibliography}{9}

\bibitem{MZ16}
E. Makai, Jr. and J. Zem\'anek,
\emph{Nice connecting paths in connected components of sets of algebraic
elements in a Banach algebra},
Czechoslovak Math. J. \textbf{66} (2016), 821--828;
arXiv:1601.01505.

\bibitem{MZ18}
E. Makai, Jr. and J. Zem\'anek,
\emph{On the structure of the set of algebraic elements in a Banach algebra
and their liftings},
arXiv:1807.01552 (2018).

\end{thebibliography}

\end{document}
